\section{Introduction}\label{sec:intro}

Time delays are inherent in a wide range of engineering systems: regenerative
machine-tool vibrations~\cite{tobias1958theory,tlusty1963stability,altintas1995analytical,munoa2016chatter},
digitally controlled and teleoperated machines, active vibration suppression,
wheel shimmy, and traffic or population dynamics all lead to delay differential
equations (DDEs)~\cite{bellman1963differential,stepan1989retarded,hale1993introduction,michiels2014stability}.
Since the characteristic function $\Dfun(\lam)$ of even the simplest linear
autonomous DDE is transcendental, it possesses infinitely many characteristic
roots, and the stability of an equilibrium can rarely be decided in closed
form. Engineering practice therefore relies on \emph{stability charts}:
maps over a plane of design parameters that show where the equilibrium is
asymptotically stable~\cite{insperger2011semidiscretization,sipahi2011stability}.
The practical value of such a chart depends on two properties of the
underlying numerical method: how fast a single parameter point can be
classified, and how much extra information (dominant root location, spectral
gap, robustness margin) the classification yields.

A large family of methods operates in the time domain by discretizing either
the solution operator or the infinitesimal generator of the delay system.
The semi-discretization method of Insperger and
St\'ep\'an~\cite{insperger2002semi,insperger2004updated,insperger2008higher,insperger2011semidiscretization},
the full-discretization method~\cite{ding2010full}, Chebyshev
collocation~\cite{butcher2004stability}, pseudospectral
differencing~\cite{breda2005pseudospectral,breda2015stability}, and the
spectral element method~\cite{khasawneh2011spectral} all construct a finite
transition matrix whose eigenvalues approximate the characteristic
multipliers or exponents; spectral methods to certify all roots in a
prescribed right half-plane are also available~\cite{wu2012reliably}.
Mature software packages implement these ideas: DDE-BIFTOOL~\cite{engelborghs2002numerical,sieber2016ddebiftool},
TRACE-DDE~\cite{breda2009tracedde}, Knut~\cite{szalai2006continuation,szalai2013knut},
and \pkg{SemiDiscretizationMethod.jl}~\cite{bachrathy2026semidisc}, which also contains
the recent multiplication-free variant with linear time
complexity~\cite{bachrathy2026multiplication}. The great strength of this
family is generality: time-periodic coefficients (e.g., milling) are handled
by the same machinery through Floquet theory. For \emph{autonomous} systems,
however, this generality is paid for at every single parameter point: a large
discretization matrix is built and its dominant eigenvalues are computed,
although only the sign of the rightmost characteristic exponent -- one bit of
information -- decides stability. Moreover, delayed systems of neutral type
(where the essential spectrum destroys the convergence of many
discretizations~\cite{hale2002strong,michiels2005eigenvalue}), systems with
algebraic constraints, i.e., delay differential--algebraic equations
(DDAEs)~\cite{campbell1980singular,ha2016analysis}, distributed
delays~\cite{cushing1977integrodifferential}, and fractional-order
systems~\cite{petras2011fractional,fioravanti2012numerical,cermak2019delay}
each require dedicated, nontrivial extensions in the time domain.

The frequency domain offers a complementary route, since all the listed
system classes share the same abstract object: an analytic characteristic
function $\Dfun(\lam)$ whose roots in the right half-plane must be counted
(for a comprehensive treatment see~\cite{gu2003stability}).
Classical tools include the D-subdivision method~\cite{neimark1949decomposition},
the analytical stability criteria of St\'ep\'an~\cite{stepan1989retarded},
the Nyquist and Mikhailov criteria~\cite{nyquist1932regeneration,mikhailov1938method},
the direct method of Walton and Marshall~\cite{walton1987direct}, the
Rekasius substitution~\cite{rekasius1980stability}, the cluster treatment
of characteristic roots~\cite{olgac2002exact}, and the frequency-sweeping
tests of Li, Niculescu and \c{C}ela~\cite{li2015analytic}; the Lambert-$W$
formalism~\cite{yi2010timedelay} and the matrix method of
Louisell~\cite{louisell2001matrix} target the dominant roots directly, the
quasi-polynomial root finder QPmR~\cite{vyhlidal2009mapping,vyhlidal2014qpmr}
computes all roots in a prescribed rectangle from the values of
$\Dfun(\lam)$ on a grid -- a contour-integral strategy also developed for
general nonlinear eigenvalue problems by Beyn~\cite{beyn2012integral} --
and the YALTA toolbox~\cite{avanessoff2013yalta} automates root-locus and
stability analysis for classical and fractional delay systems. A recent
survey of argument-principle-based counting techniques is given
in~\cite{niculescu2023counting}.

Two frequency-domain strategies are particularly relevant here. The first
solves the two real equations $\Real\Dfun(\ii\omega)=\Imag\Dfun(\ii\omega)=0$
in the parameter space extended by the frequency $\omega$, which
yields the D-curves, i.e., the parameter combinations where a root crosses
the imaginary axis. The multi-dimensional bisection method
(MDBM)~\cite{bachrathy2012bisection,bachrathy2025mdbm} traces these
one-dimensional solution manifolds very efficiently. The D-curves, however,
are \emph{not} the stability boundaries: they merely mark root crossings,
and the number of unstable roots on either side remains unknown -- and
automating that decision is surprisingly hard (triangulation into closed
cells, filtering of spurious slivers). The turning stability-lobe diagram
illustrates a more fundamental inefficiency: towards low spindle speeds the
lobe structure becomes infinitely dense, and almost all D-curves lie deep
inside the unstable domain
(cf.\ Appendix~\ref{app:gallery})~\cite{stepan2001modelling,insperger2003stability,bachrathy2013improved},
so the method spends most of its effort on curves that never bound the
stable region.

The second strategy counts the unstable roots directly by the argument
principle, i.e., by the winding number of $\Dfun(\ii\omega)$ -- the
generalized Nyquist, or Mikhailov, criterion~\cite{stepan1989retarded,kolmanovskii1999introduction}.
Root counting by contour integrals of $\Dfun'/\Dfun$ has a long history for
general analytic functions~\cite{delves1967numerical,kravanja2000computing,kravanja2000zeal};
for delay systems, Hassard~\cite{hassard1997counting} derived a
semi-analytical counting formula, Z\'{\i}tek and Vyhl\'{\i}dal formulated
argument-increment criteria for neutral systems~\cite{zitek2008argument,vyhlidal2009modification},
and, closest to the present work, Xu and Wang~\cite{xu2014exact} and Xu,
St\'ep\'an and Wang~\cite{xu2016delay} evaluated the argument-principle
``testing integral'' numerically, observing that a rough quadrature suffices
because the result must be an integer. The appeal of root counting is that
the answer $\Zint(\vect p)$ is exactly the quantity a stability chart needs.
Its curse is twofold. First, $\Zint(\vect p)$ is a piecewise-constant integer
field: it jumps abruptly across the stability boundaries, so it cannot be
interpolated, and brute-force charts need fine parameter grids. Second, the
numerical evaluation is hardest \emph{exactly at the stability boundary} --
the only region one actually cares about -- because as a root approaches the
integration line the integrand of the winding integral develops a
near-singular peak whose width shrinks to zero, and fixed grids or rough
quadratures either miss the peak or need excessive refinement; the infinite
frequency range poses an additional truncation problem.

This paper revisits the integration-based Nyquist criterion and turns it into
a fast and robust tool with \emph{interpolable} output -- fast enough that a
complete chart of a simpler system regenerates in a fraction of a second,
i.e., at near-real-time, interactive rates. The contributions are the
following. (i)~The winding integral is reformulated as a scalar ordinary
differential equation for the phase, marched by a high-order adaptive
integrator~\cite{verner2010numerically,tsitouras2011runge,rackauckas2017differentialequations}
whose step-size control resolves the near-singular peaks and crosses the
practically unbounded frequency range, with the integrand evaluated exactly
by forward-mode automatic differentiation~\cite{revels2016forward}.
(ii)~The location of the
characteristic root closest to the integration line -- normally the rightmost
root -- is estimated \emph{during the same sweep} at negligible extra cost; the
first-order estimate is the smooth, sign-correct default, and an optional local
polish (Taylor, Newton or a combined mode) sharpens it toward machine precision
for a modest fraction of the cost when an accurate root is wanted. Tracking
several local minima yields multiple roots at once. This ordered march is what
distinguishes the phase ODE from plain quadrature, whose recursive subdivision
visits the frequencies out of order and cannot recognize the minima at all.
(iii)~The root
estimate makes the chart interpolable: a combined coloring shows the integer
number of unstable roots outside and the smooth spectral gap inside the
stable domain, which is a strikingly effective visual diagnostic.
(iv)~The integrality of the computed winding number provides a built-in,
dimensionless error estimate that scales with the requested solver tolerance
and lets the accuracy of the count be prescribed in advance; where it is
structurally blind -- a skipped peak costs exactly $\pm\pi$ and leaves the
result an integer -- the independently obtained root supplies the missing
check. (v)~The estimated root distance serves
as a scalar objective for MDBM, which then traces the \emph{true} stability
boundary (not the D-curves) at very high equivalent resolution, and the most
robust parameter point is found as the center of the largest circle inscribed
in the stable domain. (vi)~The characteristic function is extracted
algorithmically from a simulation-ready right-hand side -- including singular
mass matrices (DDAEs) -- so the user never has to derive
$\Dfun(\lam)$ by hand. (vii)~The method applies verbatim to retarded,
neutral, distributed-delay, transcendental, high-dimensional and
(preliminarily) fractional-order systems, and is released as the open-source
Julia~\cite{bezanson2017julia} package \pkg{InterpolatedNyquist.jl}.
Compared with the rough-quadrature testing integral
of~\cite{xu2014exact,xu2016delay}, the present formulation adds adaptive
error control tied to the integer-residual estimate, the rightmost-root
by-product that enables interpolation and boundary tracing, and a far wider
class of demonstrated applications; compared with contour-moment methods such
as ZEAL~\cite{kravanja2000zeal}, it needs no bounded search region and is
tailored to the half-line stability sweep.

The paper is organized as follows. Section~\ref{sec:theory} summarizes the
system classes and the argument-principle counting formula.
Section~\ref{sec:method} presents the phase ODE, the rightmost-root tracking,
the higher-order refinement, and the self-validating error estimate.
Section~\ref{sec:charts} assembles these ingredients into fast stability
charts. Section~\ref{sec:extraction} describes the automatic extraction of
$\Dfun(\lam)$ from a time-domain model. Section~\ref{sec:showcase} walks
through the complete workflow on a constrained two-degree-of-freedom system
with delayed control, Section~\ref{sec:performance} quantifies accuracy, CPU
time and memory, including a comparison with the semi-discretization method,
and Section~\ref{sec:conclusion} concludes. Appendix~\ref{app:gallery} collects twelve
further case studies from all supported system classes.